\begin{table}[htbp]
\centering
\small
\caption{Mean access scores, composite index and binding modal constraint.}
\label{tab:t6_access}
\begin{tabular}{llrrrrrlr}
\toprule
City & $n$ & Major road & Rail goods shed & IWT terminal & ICD & MFAI & Binding mode & Margin \\
\midrule
Patna & 12,807 & 0.841000 & 0.776000 & 0.845000 & 0.607000 & 0.728000 & ICD & 0.169000 \\
Indore & 48,884 & 0.852000 & 0.898000 & 0.050000 & 0.702000 & 0.809000 & ICD & 0.150000 \\
Varanasi & 9,289 & 0.779000 & 0.893000 & 0.879000 & 0.792000 & 0.814000 & Major road & 0.013000 \\
Kochi & 27,759 & 0.879000 & 0.841000 & 0.830000 & 0.786000 & 0.830000 & ICD & 0.055000 \\
Kolkata & 24,852 & 0.891000 & 0.899000 & 0.907000 & 0.892000 & 0.893000 & Major road & 0.001000 \\
\bottomrule
\end{tabular}
\begin{tablenotes}\footnotesize
\item Access scores $a = T/(T+d)$, range 0–1, higher is better. MFAI is the three-mode index (road, rail, ICD), comparable across cities. The binding mode is the lowest-mean land mode; the margin is the gap to the next lowest. A margin below roughly 0.05 indicates no mode binds distinctly and the ordering should not be read as a finding. Indore's IWT score is a constant (Section 3.3).
\end{tablenotes}
\end{table}
